\chapter{Proposed Methodology}
\label{ch:method}

\section{Overview}

The proposed solution is a dedicated verification layer that sits
between an LLM--driven AutoML planner and the executable Python
pipeline that the coding agent emits. Instead of judging the
pipeline only on whether it compiles, runs or reaches a target
metric, the layer answers a strictly stronger question: does the
emitted code mechanically implement the plan it was prepared from?

The methodology is organised around five stages:

\begin{enumerate}
    \item A \textbf{task specification} that captures everything a
    plan is allowed to assume about the problem.
    \item A \textbf{Reasoning Intermediate Representation (RIR)}
    that lowers the plan into atomic typed claims.
    \item A \textbf{Program Intermediate Representation (PIR)}
    that lowers the generated code into a directed attributed
    multigraph covering AST, dataflow and library semantics.
    \item A \textbf{semantic extraction and mapping} stage that
    normalises operator calls into a bounded predicate language and
    constructs claim--to--witness mappings.
    \item A \textbf{verdict and metrics} stage that assigns
    three--valued verdicts per claim and reports coverage, weighted
    alignment, critical mismatch and localisation measures.
\end{enumerate}

The remainder of this chapter describes each stage, specifies the
alignment relation and metrics formally, presents the mismatch
taxonomy, and lists the baselines against which the verifier is
compared.

\begin{quote}
\textbf{Scope of this chapter.} This chapter specifies the
\emph{proposed} framework in full. It is a design, not a report of
delivered software. What is actually built and evaluated in this
project is a single, deliberately narrow check on the \texttt{main}
branch (Chapter~\ref{ch:impl}): for each meta--feature--guided run it
tests operation \emph{presence} (whether each action the trace names
appears in the generated code) and meta--feature \emph{value
grounding} (whether each cited statistic matches the real
meta--features within a tolerance), and returns a single boolean
faithfulness verdict. Everything else described below --- the
graph--based (multigraph) Program IR, the set--valued mapper, the
three-- and four--valued verdicts, the full mismatch taxonomy of
Table~\ref{tab:taxonomy}, the evaluation metrics and the baseline
comparison --- is proposed and is \emph{not} implemented or evaluated
in this project. The reader should treat every mechanism below as
proposed unless Chapter~\ref{ch:impl} marks it implemented.
\end{quote}

\section{System Requirements}

\subsection{Hardware Requirements}

The verifier itself is CPU--only and is comfortable on commodity
developer hardware. The LLM backends used to generate plans and code
are the main hardware driver.

\begin{itemize}
    \item A 64--bit workstation or server with at least 16~GB of RAM
    and an SSD for pipeline execution and corpus storage.
    \item An LLM backend. The reported experiments used hosted
    \texttt{Ollama} cloud models (for example a 120B model), so no
    local GPU was required; running a model locally instead needs a
    CUDA--capable GPU with enough memory for that model.
    \item Optional access to a Redis--capable host for job
    orchestration when the pipeline is deployed as a service.
\end{itemize}

\subsection{Software Requirements}

\begin{itemize}
    \item Python 3.10 or later, with
    \texttt{scikit--learn}~\cite{sklearn}, \texttt{pandas},
    \texttt{numpy} and \texttt{pydantic}. The implemented verifier
    uses only the standard--library \texttt{ast} module; the proposed
    graph--based Program IR additionally requires \texttt{networkx}
    and \texttt{libcst}, which are not yet a project dependency.
    \item \texttt{FastAPI} (with \texttt{uvicorn}) and
    \texttt{rq}/\texttt{Redis} for the HTTP services, job queueing and
    run--worker.
    \item An LLM runner with a configurable backend (the reported
    runs used hosted \texttt{Ollama} cloud models) so that the
    generation model can be swapped without changing the rest of the
    pipeline.
    \item A bounded execution environment for running generated code
    (a separate Python subprocess with a wall--clock timeout).
    Container--level isolation with resource and filesystem limits is
    proposed, not implemented.
    \item LaTeX, for reproducing this report and generated figures.
\end{itemize}

\section{Proposed Approach}

The verifier consumes four objects: a \emph{task specification} $T$,
a \emph{reasoning artifact} $R$, a \emph{code artifact} $C$ and an
\emph{execution trace} $E$. It produces a per--claim verdict table,
a mismatch summary and a set of aggregate metrics.

\subsection{Task Specification}

The input records the dataset reference, the target column, the task
type (binary / multiclass / regression), the allowed or prohibited
models, the preprocessing constraints, the evaluation metric, the
split policy, the resource limits and the required outputs. Each
field receives a provenance tag --- \texttt{USER},
\texttt{INFERRED} or \texttt{DEFAULT} --- so that an unstated agent
assumption is never silently treated as a user requirement.

\subsection{Reasoning Intermediate Representation (RIR)}

Unrestricted prose cannot be used as a mechanical specification. The
planner (or an extraction step applied to free--form reasoning)
emits a readable rationale together with a structured RIR. Each
atomic claim $r_i$ has the form
\[
r_i = \langle id, o, s, a, \theta, d, q, w, \pi \rangle ,
\]
where $o$ is an operation class, $s$ is a scope, $a$ is an operator
family, $\theta$ are parameters, $d$ is an ordering or dependency
predicate, $q$ is a qualifier, $w$ is a criticality and $\pi$ is
provenance. For example, a claim can be written as
\begin{verbatim}
R2: TRANSFORM(scope=numeric,
              family=standardization,
              fit=train_only,
              before=model, weight=2).
\end{verbatim}
Claims cover loading, splitting, imputation, encoding, scaling,
feature selection, dimensionality reduction, model choice, fitting,
prediction, evaluation and constraints. Meaningful but
non--operational claims are retained as \texttt{NONEXECUTABLE} and
are not counted as aligned.

Free--form plans require a separate extraction step. To avoid
circular evaluation, extraction errors are measured against
human--annotated claims; direct structured emission is the primary
condition and free--form extraction is a harder comparison.

\paragraph{What the reasoning trace actually is (and its limitation).}
The design above treats the reasoning as a plan produced \emph{before}
the code. In the system actually built and evaluated here, the
reasoning trace is not captured independently of the code. Under the
meta--feature--guided (B2) condition the LLM is asked to return a
single, self--contained Python script that, when run, prints a
\texttt{SCORE:} line and then a \texttt{REASONING:} line carrying a
small JSON object of decisions (the step, the cited meta--feature, the
\texttt{observed\_value} and the \texttt{action}). The faithfulness
check reads that JSON from the script's output and tests it against the
same script's code and the real meta--features. The trace is therefore
\emph{post hoc}: the code and its stated reasoning come from one
generation step, so the check measures the internal consistency of a
single LLM response --- whether the script cites an action it did not
import, or a number it did not read correctly --- rather than the
divergence of code from a plan fixed in advance by a human or by a
separate planning stage. Capturing an independent pre--code plan would
require a distinct planning step and is left to future work.

\subsection{Code Generation and Program IR}

The coding agent receives the task specification and the plan, and
emits source code without verifier labels. Source, environment, seed
and transcript are versioned, and failed programs are preserved as
outcomes.

Python is parsed into an AST and lowered into a directed attributed
multigraph
\[
\mathcal{P} = (V_P, E_D, E_C, \lambda, \eta),
\]
where $V_P$ contains data objects, transforms, estimators, splits,
metrics and predicates; $E_D$ records dataflow; $E_C$ records
control or calls; $\lambda$ maps nodes to normalised semantics; and
$\eta$ stores arguments, locations and types.

An AST identifies syntax but inadequately represents pipeline order.
A runtime trace observes only one path. The combined representation
uses static analysis for coverage and optional instrumentation to
resolve dynamic dispatch, fitted columns and values. Library
adapters normalise \texttt{Pipeline},
\texttt{ColumnTransformer}, estimators, scoring strings and custom
transformers.

\section{Algorithms / Models}

\subsection{Semantic Extraction, Mapping and Verdicts}

Calls are normalised into predicates such as
\begin{verbatim}
    IMPUTE(num, strategy=median),
    SCALE(num, semantics=zscore, fit=train),
    MODEL(family=logistic, balanced=true),
    EVALUATE(metric=f1, average=macro).
\end{verbatim}
Bodies or trusted summaries resolve aliases and wrappers. Constant
propagation is conservative; unknown values are left symbolic.
Dataflow provenance enables leakage checks --- for example, whether a
scaler was fitted before the train/validation split.

For each claim, the mapper constructs program witnesses from the
operation type, scope, dataflow neighbourhood and semantic family.
The rules then check presence, operator semantics, parameters, order,
dependency, target and metric constraints. Each report records the
claim, the expected predicate, the observed evidence, the source span,
the verdict, the category and the severity.

Verdicts are reported at two levels. Each individual claim receives a
\emph{three--valued} verdict --- \textsc{aligned}, \textsc{misaligned}
or \textsc{unverifiable} --- the last when there is not enough
evidence to decide it. The claim verdicts are then aggregated, with
criticality gating, into a \emph{four--valued} pipeline verdict:
\begin{itemize}
    \item \textsc{Aligned} --- all decided claims align and no
    contradiction is present.
    \item \textsc{Partially aligned} --- only non--critical claims
    mismatch or remain unaccounted for.
    \item \textsc{Misaligned} --- a critical claim fails or a
    contradiction is present.
    \item \textsc{Unverifiable} --- there is not enough evidence to
    decide the pipeline either way.
\end{itemize}
Abstinence (\textsc{Unverifiable}) prevents uncertainty from turning
into false certainty.

\subsection{Alignment Relation}

Let $\mathcal{R} = \{r_1, \dots, r_n\}$ be the set of plan claims and
$V_P = \{p_1, \dots, p_m\}$ the set of program nodes. The mapper is
set--valued:
\[
M : \mathcal{R} \to 2^{V_P} \cup \{\emptyset\}.
\]
A claim can require several nodes; leakage--safe scaling, for
example, requires a split, a scaler fitted on the training branch and
reuse on validation data.

The per--claim relation is three--valued,
\[
A(r_i, M(r_i)) \in \{1, 0, \bot\},
\]
where $1$ denotes agreement, $0$ an established mismatch and $\bot$
insufficient evidence. It combines the applicable predicates,
\[
A_i = A_i^{\mathrm{pres}} \wedge A_i^{\mathrm{sem}} \wedge
A_i^{\mathrm{param}} \wedge A_i^{\mathrm{ord}} \wedge
A_i^{\mathrm{dep}} .
\]
A three--valued conjunction returns $0$ if any required predicate is
false, $1$ if all are true, and $\bot$ otherwise. In the implemented
\texttt{main}--branch check only the first two conjuncts are evaluated
--- presence ($A^{\mathrm{pres}}$) and meta--feature value grounding
(a form of $A^{\mathrm{sem}}$) --- and they are reduced to a single
boolean per run; the parameter, ordering and dependency conjuncts
belong to the proposed design.

\section{Evaluation Metrics}

\emph{These metrics are part of the proposed design. None of them is
computed in this project's experiment; Chapter~\ref{ch:results} reports
success rate, predictive score, error categories and the boolean
faithfulness verdict, not the quantities below.} They are defined here
so that a fuller verifier, if built, would have a precise and
non--inflatable reporting target.

Let $v_i = 1$ if claim $i$ is decided, $a_i = 1$ only for a decided
match, and $w_i > 0$ be a criticality weight assigned per claim type
(for example, a leakage--relevant ordering claim carries more weight
than an optional feature--selection claim).

\begin{itemize}
    \item \textbf{Claim coverage} --- the fraction of plan claims the
    verifier was able to decide, rather than leave unverifiable:
    \[
    C_{\mathrm{claim}} = \frac{1}{n}\sum_{i=1}^{n} v_i.
    \]
    \item \textbf{Weighted semantic alignment} --- among the decided
    claims, the share that matched, with more important claims (for
    example a leakage--relevant ordering claim) weighted more heavily:
    \[
    S_{\mathrm{align}} = \frac{\sum_i w_i v_i a_i}{\sum_i w_i v_i}
    \]
    when the denominator is non--zero. Coverage must be reported
    alongside the score.
    \item \textbf{Critical mismatch rate} --- the fraction of
    safety--critical claims (the set $K$) that were contradicted:
    \[
    \mathrm{CMR} = \frac{\sum_{i\in K}\mathbf{1}[A_i = 0]}{|K|}.
    \]
    \item \textbf{Localisation} --- how well predicted mismatch
    categories and their source locations match human labels, measured
    by category macro--F1 and source--span overlap. False--alignment
    and false--misalignment rates are reported separately.
\end{itemize}

Coverage is reported alongside the alignment score so that a high
value cannot be achieved by labelling most claims
\textsc{unverifiable} and declaring the handful of remaining claims
aligned.

\subsection{Semantic Equivalence}

Alignment cannot be decided by string equality. The design uses four
levels of evidence: \emph{lexical} (which symbols are candidate
matches), \emph{structural} (AST and graph patterns, or an operator's
role in the program), \emph{semantic} (the operator's meaning, taken
from a trusted summary or an equivalent rewrite), and
\emph{behavioural} (agreement on a bounded set of probe inputs over a
declared input domain). Lexical evidence alone cannot certify a claim,
and behavioural agreement over probes does not prove universal
semantic equivalence. Without a trusted contract, a custom function
that relies on reflection, external state or undecidable dispatch is
labelled $\bot$ (unverifiable).

For instance, \texttt{StandardScaler} and a custom $(x - \mu)/\sigma$
may match if such code fits $\mu$ and $\sigma$ on training data,
uses the same variance and zero--variance conventions, and transforms
only the intended columns. Addition and subtraction observed on their
own are not enough.

\subsection{Worked Example}
\label{sec:worked-example}

To make the representations concrete, consider a single
preprocessing decision on the Titanic dataset. The LLM is given the
dataset meta--features and emits the natural--language reasoning:

\begin{quote}\itshape
``The \texttt{Fare} column is heavily skewed with many outliers
(outlier ratio $0.128 > 0.05$), so I scale it with
\texttt{RobustScaler} rather than \texttt{StandardScaler}.''
\end{quote}

\noindent\textbf{(1) RIR claim.} This lifts to one atomic typed
claim:
\begin{verbatim}
C1: TRANSFORM(op=scale, scope={Fare},
              family=robust_scaling,
              meta_feature=
                distributional.outlier_ratio_per_numeric.Fare,
              observed_value=0.128, threshold=0.05,
              action=RobustScaler, weight=medium).
\end{verbatim}

\noindent\textbf{(2) Generated code.} The coding step emits, among
other lines:
\begin{lstlisting}[language=Python]
from sklearn.preprocessing import RobustScaler
from sklearn.compose import ColumnTransformer
pre = ColumnTransformer([("num", RobustScaler(), ["Fare"])])
\end{lstlisting}

\noindent\textbf{(3) PIR.} The program lowers to nodes including a
call \texttt{RobustScaler()} imported from
\texttt{sklearn.preprocessing}, wired into a
\texttt{ColumnTransformer} whose column literal set is
\texttt{\{Fare\}}.

\noindent\textbf{(4) Mapping and verdict.} The claim \texttt{C1} maps
to the \texttt{RobustScaler} node. The verifier checks: the cited
action \texttt{RobustScaler} is present in the code
($A^{\mathrm{pres}}=1$); the cited
\texttt{observed\_value}~$0.128$ equals the real meta--feature within
tolerance ($A^{\mathrm{sem}}$ value--grounding~$=1$); and, in the
full design, the applied column \texttt{Fare} and the fit--after--split
ordering also hold. The verdict is \textsc{aligned}.

\noindent\textbf{(5) The misalignment case.} Had the code instead
used \texttt{StandardScaler} while the trace still claimed
\texttt{RobustScaler}, then $A^{\mathrm{pres}}=0$ and the verdict is
\textsc{misaligned} --- even though the pipeline compiles, runs and
may reach an identical score. This is precisely the failure that
execution-- and performance--based checks cannot see, and it is the
check that the implemented \texttt{main}--branch verifier
(Chapter~\ref{ch:impl}) already performs. Figure~\ref{fig:rirpir}
illustrates the full mapping.

\begin{figure}[H]
    \centering
    \miroplaceholder{figures/project_figures/rir_pir_mapping.png}{%
    RIR--to--PIR mapping for the worked example: the natural--language
    reasoning on the left, its atomic RIR claim, the generated code
    and its PIR (AST call node + column literals + import) on the
    right, with the claim--to--witness arrows and the resulting
    per--check verdicts (presence / value / column / order).}
    \caption{Worked RIR--PIR claim--to--code mapping for a single
    \texttt{RobustScaler} decision on the \texttt{Fare} column.}
    \label{fig:rirpir}
\end{figure}

\section{Mismatch Taxonomy}

The proposed taxonomy of reasoning--code misalignment is given in
Table~\ref{tab:taxonomy}. It is defined so that the mechanical
condition of each category can be checked from RIR and PIR evidence,
and so that each category can be weighted by severity. Secondary
tags are permitted; placing a scaler before a split is both an
ordering and a data--handling mismatch.

\begin{table}[H]
\centering
\caption{Proposed taxonomy of AutoML reasoning--code misalignment.}
\label{tab:taxonomy}
\small
\renewcommand{\arraystretch}{1.2}
\begin{tabular}{p{3.3cm} p{4cm} p{4.5cm} p{2.2cm}}
\toprule
\textbf{Category} & \textbf{Mechanical condition} &
\textbf{Example} & \textbf{Severity} \\
\midrule
Missing implementation & No reachable witness &
Feature selection planned but absent & Medium--High \\
Contradictory implementation & Observed predicate conflicts &
Preserve rows planned; code drops missing rows & High \\
Model mismatch & Estimator family incompatible &
Logistic model planned; forest used & High \\
Parameter mismatch & Operator matches but value differs &
Five folds stated; code uses three & Medium--High \\
Pipeline--order mismatch & Precedence or fit scope violated &
Scaling fitted before split & High \\
Evaluation mismatch & Metric differs &
Macro F1 stated; accuracy used & High \\
Data--handling mismatch & Columns, target, split or sampling differs
& Stratification absent & High \\
Constraint mismatch & Resource or admissibility rule violated &
Sparse model required; ensemble used & Medium--High \\
Unsupported claim & No operational definition &
Use a stable model & Unverifiable \\
Dead reasoning & Matching node is off effective path &
Selector constructed but unused & Medium \\
Extraneous operation & Consequential unstated step added &
Target--dependent preprocessing added & Contextual \\
\bottomrule
\end{tabular}
\end{table}

\section{Baselines}

The proposed evaluation compares five systems. This comparison is part
of the proposed protocol and is \emph{not} carried out in this project
(see Chapter~\ref{ch:results}); it is described here to define what a
fair assessment of the fuller verifier would involve.

\begin{enumerate}
    \item \textbf{Execution--only}: aligned if execution succeeds.
    \item \textbf{Performance--only}: aligned if execution completes
    and the requested performance threshold is met.
    \item \textbf{LLM judge}: a separate model assesses alignment by
    reasoning over the plan and code without access to program
    analysis tools.
    \item \textbf{Keyword/API matching}: claims are considered
    aligned when they match Python method names and metric strings.
    \item \textbf{Mechanical verifier}: performs full RIR--PIR
    mapping with semantic normalisation, program--graph constraints,
    and uncertainty estimation.
\end{enumerate}

Table~\ref{tab:baselines} summarises what each baseline is capable of
doing.

\begin{table}[H]
\centering
\caption{Baseline capabilities.}
\label{tab:baselines}
\renewcommand{\arraystretch}{1.2}
\begin{tabular}{lcccc}
\toprule
\textbf{Method} & \textbf{Runs} & \textbf{Plan} & \textbf{Graph} &
\textbf{Localizes} \\
\midrule
Execution--only & \checkmark & -- & -- & -- \\
Performance--only & \checkmark & Part. & -- & -- \\
LLM judge & Opt. & \checkmark & -- & Text \\
Keyword match & -- & \checkmark & -- & Token \\
Mechanical verifier & Opt. & \checkmark & \checkmark & Claim \\
\bottomrule
\end{tabular}
\end{table}

\section{Data Collection and Preprocessing}

Two corpora are needed.

\textbf{Natural corpus.} Unaltered LLM plans and code. Claims are
extracted and program evidence is mapped by two annotators
independently; a third has the final say when disagreements occur.
Each annotator estimates alignment, category, severity and source
span. Reported numbers include (1) agreement per claim extraction,
(2) ternary labels and (3) taxonomy classes.

\textbf{Seeded--mutation corpus.} Well--aligned pipelines are the
basis. A controlled mutation deletes an operator, swaps a model or a
metric, changes a parameter, factors out transforms, alters the
split policy, disconnects one step from another and violates any
constraints. Mutations maintain syntax, and some subset must
guarantee successful execution. This is a direct test of the
motivating case. Negative controls consist of equivalent
refactorings and aliases, wrappers, and semantically equivalent
implementations. Mutation templates must retain names and comments
(if possible) to prevent keyword baselines from taking advantage of
simple cues; the point is that when code is mutated, only a semantic
change is introduced where humans intended, and nothing else.

\section{Summary}

The proposed methodology formalises reasoning--code alignment as a
set--valued mapping from atomic typed claims (RIR) to program
witnesses in a directed attributed multigraph (PIR), with a
three--valued verdict space that includes a dedicated
\textsc{unverifiable} escape. A bounded predicate language
normalises scikit--learn calls, an AutoML--specific taxonomy
describes mismatch categories, and a \emph{proposed} metric suite
(coverage, weighted alignment, critical mismatch and localisation)
would summarise the result in a way that cannot be inflated by
abstention. A set of baselines, including an LLM--judge, is specified
so that the added value of the mechanical check could be measured
rather than assumed. Of all this, only the presence and
value--grounding check is implemented and evaluated in this project;
the next chapter presents the system architecture within which that
check runs on \texttt{main}, while the fuller design remains proposed.
